\section{Related Literature}
\label{sec:related-literature}

\subsection{Defaults and choice architecture}
That pre-set options move consequential choices is among the most
replicated results in economics, from retirement saving
\parencite{madrian2001power, carroll2009optimal} and organ donation
\parencite{johnson2003defaults} to the broad program of
\textcite{thaler2008nudge}; \textcite{dellavigna2009psychology}
surveys the field. \textcite{haggag2014default} brought the result to
tipping using the same New York taxi system we study, and
\textcite{alexander2021recommendations} showed in the same setting
that raising the suggested-tip menu raises tips---about seventeen
cents per percentage point of suggestion---at the cost of more zero
tips. Our contribution is orthogonal to the menu \emph{level}: we hold
the menu fixed and study the \emph{base} it multiplies.

\textcite{thakral2026five} are closest to us and should be read
alongside this paper. Using the September 2012 fare change and the two
vendors' differing treatment of tax and tolls, they establish that a
larger base raises tips even among riders who do not press a default
button, and that add-ons crowd tips in or out according to whether the
interface includes them---facts our vendor evidence in
Section~\ref{sec:results-vendor} and Appendix~\ref{app:replication}
reproduce. Their question is what tipping \emph{is}: they build a model
of a default-anchored norm with a convex pain of paying. Ours is what a
\emph{tax} does: their identifying variation is a fare change and a
route-determined toll, while ours is a legislated surcharge that
switches at a known minute, at a dose the legislature set and then
changed by 150 percent. That difference is what lets us report a
pass-through rate per statutory dollar, test its stability across doses
and directions, extend it to two congestion-pricing programs, and
aggregate the transfer. Choice architecture here is not an anchor but a
transmission mechanism from fiscal policy to a private transfer.

\subsection{Tax salience, attention, and shrouding}
\textcite{chetty2009salience} and \textcite{finkelstein2009ezpass}
established that how a tax is displayed changes how agents respond to
it; \textcite{taubinsky2018attention} show attention to sales taxes is
endogenous and heterogeneous, \textcite{feldman2015tax} and
\textcite{goldin2013smoke} develop the toolkit for partially attentive
consumers, and \textcite{morrison2023rules} document rule-of-thumb
responses to price changes. \textcite{gabaix2006shrouded} supply the
market logic of add-on prices that firms do not compete away. Our
channel is distinct from all of these in one respect: the surcharge is
fully disclosed, itemized on the screen, and riders still tip on
it---not because the tax is hidden, but because the interface has
already folded it into the arithmetic of an unrelated voluntary
payment. It is shrouding by \emph{composition} rather than by
concealment, and the December 2022 test shows riders undo only a fifth
of a 150 percent surcharge increase.

\subsection{Tipping}
Tipping is a large voluntary transfer whose variation service quality
explains only partly \parencite{lynn2000tipping, azar2020tipping};
\textcite{conlin2003norm} model it as norm enforcement, and
\textcite{chandar2019drivers}, in Uber's national rollout, find rider
fixed traits dominate trip attributes---most riders never tip, a few
always do. Field evidence on point-of-sale prompts
\parencite{warren2021manipulated} documents consumer irritation with
suggestion screens. The channel is not specific to taxis:
\textcite{lynn2023largesample}, in 68 million restaurant checks, find
tipping rates rise with the local sales-tax rate---which is what tipping
on a tax-inclusive check total implies, in an industry with no taxi
meter in sight. Against this backdrop our estimates show that a
substantial share of observed tip variation is set by payment-terminal
arithmetic: a dollar of government surcharge moves the average tip by
fifteen cents regardless of which vendor's screen, which dose, or
which direction.

\subsection{Congestion pricing}
Evaluations of cordon pricing measure traffic, revenue, and
environmental margins, from London \parencite{leape2006london} to the
January 2025 New York program, whose short-run traffic effects
\textcite{cook2025congestion} measure with network speed data. Fiscal
incidence in these evaluations ends at the statutory payment. We
document a channel none of them counts: the toll enters the tip base,
so part of the rider's payment continues past the MTA to the driver.
Roughly \$2.2 million a year of New York's congestion toll becomes
gratuities---small beside program revenue, but a transfer created
entirely by interface design, and one that generalizes to every
percentage-suggestion payment system that sits downstream of a
government fee.
